\chapter{ภาคผนวก —การแยกเชื้อและทดสอบจุลินทรีย์ทางการเกษตร}
\label{ch:appendix_lab}

ภาคผนวกนี้รวบรวมขั้นตอนปฏิบัติการทดลองพื้นฐานที่จำเป็นสำหรับการแยก การคัดเลือก และการทดสอบประสิทธิภาพของจุลินทรีย์ที่มีบทบาทสำคัญทางการเกษตร เพื่อเป็นคู่มือสำหรับนักศึกษาในการทำวิจัยและการเรียนการสอนภาคปฏิบัติการ

\section{ปฏิบัติการที่ 1 —การแยกเชื้อ Rhizobium จากปมรากถั่ว}

แบคทีเรียสกุล \organism{Rhizobium} อาศัยอยู่ในปมรากพืชตระกูลถั่ว ในการแยกเชื้อจำเป็นต้องทำความสะอาดผิวปมรากเพื่อกำจัดแบคทีเรียดินชนิดอื่นภายนอกก่อนบดและแยกเชื้อ

\begin{labprotocol}[ขั้นตอนการแยกเชื้อ Rhizobium จากปมรากถั่ว]
\textbf{วัสดุและอุปกรณ์} ปมรากถั่วสดที่มีสีชมพู (แสดงถึงการมีเลกเฮโมโกลบินที่ทำงานอยู่), แอลกอฮอล์ 70\%, น้ำยาฟอกขาว (sodium hypochlorite) 1--2\%, น้ำกลั่นนึ่งฆ่าเชื้อ, อาหารเลี้ยงเชื้อ Yeast Extract Mannitol Agar (YEMA) ผสมสีย้อมคอนโกเรด (Congo red), ปากคีบ, แท่งแก้วบดละเอียดนึ่งฆ่าเชื้อ

\textbf{ขั้นตอนการทดลอง}
\begin{enumerate}
    \item ล้างปมรากถั่วด้วยน้ำสะอาดเพื่อกำจัดเศษดินออกให้หมด เลือกปมที่สมบูรณ์ไม่ฉีกขาด
    \item จุ่มปมรากในแอลกอฮอล์ 70\% เป็นเวลา 10 วินาที เพื่อลดแรงตึงผิวและฆ่าเชื้อเบื้องต้น
    \item ย้ายปมรากไปแช่ในน้ำยาฟอกขาว 1--2\% นาน 3--5 นาที เพื่อฆ่าเชื้อที่ผิว (surface sterilization)
    \item ล้างปมรากด้วยน้ำกลั่นนึ่งฆ่าเชื้ออย่างน้อย 5 ครั้ง เพื่อล้างน้ำยาฟอกขาวที่ตกค้างออกให้หมด
    \item ใช้ปากคีบคีบปมรากใส่ในถ้วยแก้วจิ๋วหรือจานเพาะเชื้อนึ่งฆ่าเชื้อ เติมน้ำกลั่นนึ่งฆ่าเชื้อ 2--3 หยด บดปมรากให้ละเอียดด้วยแท่งแก้วบด
    \item ใช้ลูป (inoculating loop) แตะน้ำของเหลวที่ได้จากการบด นำไปขีดบนจานอาหาร YEMA ผสมสีย้อมคอนโกเรด
    \item บ่มที่อุณหภูมิ 28--30\degree C เป็นเวลา 3--5 วัน สังเกตโคโลนีของ \organism{Rhizobium} ซึ่งจะมีลักษณะกลม นูน เมือกขุ่นหรือกึ่งใส และมีสีขาวหรือสีชมพูอ่อนมาก (ไม่ดูดซับสีย้อมคอนโกเรดสีแดงเข้มเหมือนแบคทีเรียปนเปื้อนชนิดอื่น)
\end{enumerate}
\end{labprotocol}

\section{ปฏิบัติการที่ 2 —การคัดเลือกจุลินทรีย์ละลายฟอสเฟตบนอาหาร Pikovskaya}

การทดสอบความสามารถในการละลายฟอสเฟตทำได้โดยใช้สารประกอบแคลเซียมฟอสเฟตที่ละลายน้ำยาก ผสมในอาหารเลี้ยงเชื้อจนทำให้เกิดสีขาวขุ่น หากจุลินทรีย์ละลายฟอสเฟตได้ จะเกิดวงใสรอบโคโลนี (clear zone)

\begin{labprotocol}[การคัดเลือกและคำนวณดัชนีการละลายฟอสเฟต (PSI)]
\textbf{วัสดุและอุปกรณ์} ตัวอย่างเชื้อแบคทีเรียหรือราที่ต้องการทดสอบ, อาหารแข็ง Pikovskaya's medium (มีสารแคลเซียมฟอสเฟต $Ca_3(PO_4)_2$ ขุ่นผสมอยู่), ไม้บรรทัดหรือเวอร์เนียคาลิปเปอร์

\textbf{ขั้นตอนการทดลอง}
\begin{enumerate}
    \item จุดหรือเพลทเชื้อตัวอย่างลงบนจานอาหาร Pikovskaya เป็นจุดวงกลม (spot inoculation)
    \item บ่มที่อุณหภูมิ 30\degree C เป็นเวลา 5 วัน (กรณีแบคทีเรีย) หรือ 7 วัน (กรณีเชื้อรา)
    \item สังเกตการเจริญเติบโตและการเกิดวงใสรอบโคโลนีเนื่องจากการหลั่งกรดอินทรีย์ไปทำปฏิกิริยากับแคลเซียมฟอสเฟต
    \item วัดขนาดเส้นผ่านศูนย์กลางของโคโลนี ($d_c$) และเส้นผ่านศูนย์กลางของวงใสรวมโคโลนี ($d_z$) จากนั้นนำมาคำนวณดัชนีละลายฟอสเฟต (Phosphate Solubilization Index —PSI)
\end{enumerate}

\begin{equation}
\text{PSI} = \frac{\text{เส้นผ่านศูนย์กลางของวงใส (} d_z \text{)}}{\text{เส้นผ่านศูนย์กลางของโคโลนี (} d_c \text{)}}
\end{equation}
\end{labprotocol}

\begin{worked}[ตัวอย่างที่ 10.1 การคำนวณดัชนี PSI]
แบคทีเรีย \organism{Bacillus subtilis} สายพันธุ์ A เจริญเติบโตบนจานอาหาร Pikovskaya นาน 5 วัน วัดขนาดโคโลนีได้ 8 มิลลิเมตร และวัดขนาดวงใสรอบโคโลนีได้กว้างรวม 16 มิลลิเมตร จะได้ค่าดัชนีความสามารถในการละลายฟอสเฟต (PSI) เท่ากับ
$$\text{PSI} = \frac{16\text{ mm}}{8\text{ mm}} = 2.0$$
สายพันธุ์ที่มีค่า PSI มากกว่า 2.0 ถือเป็นสายพันธุ์ที่มีประสิทธิภาพสูง เหมาะต่อการนำไปพัฒนาต่อเป็นปุ๋ยชีวภาพ
\end{worked}

\section{ปฏิบัติการที่ 3 —การทดสอบกิจกรรมปฏิปักษ์ของ Trichoderma ด้วยเทคนิค Dual Culture}

เทคนิค Dual culture ใช้สำหรับประเมินความสามารถของเชื้อปฏิปักษ์ เช่น \organism{Trichoderma} ในการเจริญเติบโตเข้าครอบครองและแข่งขันกับเชื้อราโรคพืชเป้าหมาย เช่น \organism{Phytophthora} หรือ \organism{Sclerotium}

\begin{labprotocol}[ขั้นตอนการทำ Dual Culture Technique]
\textbf{วัสดุและอุปกรณ์} อาหาร Potato Dextrose Agar (PDA), เชื้อราปฏิปักษ์ \organism{Trichoderma}, เชื้อราก่อโรคพืชตัวแทน (เช่น \organism{Sclerotium rolfsii}), ที่เจาะเนื้อเยื่อวุ้นนึ่งฆ่าเชื้อ (cork borer) ขนาด 5 มิลลิเมตร

\textbf{ขั้นตอนการทดลอง}
\begin{enumerate}
    \item ใช้ cork borer เจาะที่ขอบโคโลนีราปฏิปักษ์ และขอบโคโลนีราก่อโรคพืชที่เจริญดีแล้ว
    \item นำแผ่นวุ้นเชื้อทั้งสองมาวางบนจานอาหาร PDA จานเดียวกัน โดยให้ห่างกันระยะประมาณ 5 เซนติเมตร (แต่ละจุดห่างจากขอบจานเพาะเชื้อ 1.5 เซนติเมตร)
    \item ทำชุดควบคุมโดยวางแผ่นวุ้นเชื้อราก่อโรคพืชเพียงอย่างเดียวที่ตำแหน่งเดิม
    \item บ่มที่อุณหภูมิห้อง สังเกตและวัดขนาดรัศมีการเจริญของเชื้อราโรคพืชทุกวันจนกระทั่งในจานควบคุม เชื้อราโรคพืชเจริญเต็มจาน
    \item คำนวณร้อยละการยับยั้งการเจริญเติบโต (Percentage of Inhibition of Radial Growth —PIRG) จากสูตร
\end{enumerate}

\begin{equation}
\text{PIRG (\%)} = \frac{R_1 - R_2}{R_1} \times 100
\end{equation}
โดยที่ $R_1$ คือ รัศมีการเจริญของเชื้อราโรคพืชในจานควบคุม และ $R_2$ คือ รัศมีการเจริญของเชื้อราโรคพืชในจานทดสอบที่มี \organism{Trichoderma} อยู่ร่วม
\end{labprotocol}

\section*{สรุปภาคผนวก}

กระบวนการแยกและทดสอบจุลินทรีย์เหล่านี้เป็นการทดสอบความเข้ากันได้ขั้นพื้นฐานทางสรีรวิทยา จุลินทรีย์ที่ผ่านการคัดกรองในระดับห้องปฏิบัติการ (in vitro) จะถูกนำไปทดสอบประสิทธิภาพการกระตุ้นพืชต่อในระดับเรือนทดลอง (greenhouse) และแปลงปลูกจริง (field trial) เพื่อเป็นผลิตภัณฑ์ชีวภัณฑ์ที่มีความน่าเชื่อถือสูง
